% !TEX root = ./main.tex
% Summary of David A. Lake (2011), "Why \"isms\" Are Evil"
% This file is intentionally a LaTeX fragment: it does not load a document class
% or packages, so it can be included from your existing main.tex/sty setup.

\section{Why "isms" Are Evil David Lake}
\cite{lakeWhyIsmsAre2011}



\subsection{Central Argument and
Purpose}

Lake's essay is a critique of how international studies is organized
intellectually and professionally. His central claim is not that
theoretical diversity is undesirable, but that scholars have turned
theoretical traditions and epistemological preferences into academic
``sects.'' Instead of using theories as limited tools for understanding
specific problems, researchers often treat broad traditions such as
realism, liberalism, constructivism, Marxism, neoliberal
institutionalism, feminism, or the English School as competing
intellectual religions. Each group tends to defend its own assumptions,
produce research that confirms those assumptions, and debate rival
traditions at the level of first principles. In Lake's view, this
sectarian organization generates less understanding rather than more. It
also distracts scholars from their fundamental social responsibility:
improving knowledge of world politics and, where possible, identifying
ways to promote more humane outcomes.

The problem is especially serious because the discipline does not
possess a single, universal and empirically powerful theory of
international politics. World politics is too complex, and existing
knowledge is too incomplete, for one general framework to explain
everything. Lake therefore urges intellectual humility. Scholars should
accept that most theories are partial and contingent, specify the
conditions under which they apply, and concentrate on mid-level
explanations of concrete phenomena. He favors ``analytic eclecticism'':
researchers should be willing to draw on different assumptions and
theories when different problems require different tools. At the
epistemological level, he makes a parallel argument. The divide between
nomological and narrative forms of explanation is real and enduring, but
neither approach is universally superior. Both can produce valid
knowledge, and the relevant question is which approach is most
illuminating for a particular research problem.

\subsection{Theory: How Research Traditions Become Academic
Sects}

International studies has long been divided into what are often called
paradigms, but Lake prefers the term ``research traditions.'' Each
tradition contains a set of core assumptions concerning the relevant
actors in world politics, their interests, the processes through which
they make decisions, and the nature of political interaction. Different
traditions may focus on individuals, groups, states, or other
collectivities; assume that actors seek wealth, power, security, status,
identity, or justice; and portray international politics as
fundamentally conflictual, cooperative, or socially constructed. These
assumptions are useful because they simplify a complex reality and help
researchers identify puzzles. The difficulty arises when scholars forget
that they are simplifications and begin to treat them as complete
theories or even as truths about the world.

Lake argues that the coexistence of many traditions should be expected.
International politics is an exceptionally complicated social system, so
it is unsurprising that different approaches illuminate different parts
of it. The discipline's weakness lies not in pluralism itself but in the
professional practices that transform pluralism into sectarianism. He
identifies five linked pathologies. Together they encourage scholars to
defend approaches rather than explain phenomena, reward intellectual
extremism, and create self-confirming research communities.

\subsection{The first pathology: reifying research
traditions}

The first pathology is the reification of research traditions. Scholars
routinely organize literature reviews, courses, graduate seminars,
handbooks, and disciplinary histories around familiar ``isms.'' This has
practical advantages because it provides a convenient map of a large
literature, but it also creates artificial boundaries. Complex
scholarship is reduced to a few recognizable schools, subtle differences
within traditions disappear, and authors are classified according to
labels they may not accept themselves. Scholars may see their own work
as eclectic and nuanced while categorizing opponents as ``realists,''
``constructivists,'' or ``liberals.'' As a result, intellectual
identities are partly imposed by others. The discipline collectively
produces the very categories that then appear to be natural and
independent objects.

\subsection{The second pathology: rewarding
extremism}

Once research traditions are reified, the discipline rewards the
clearest and most extreme statements of each tradition. Canonical works
such as Waltz's Theory of International Politics, Keohane's After
Hegemony, or Wendt's Social Theory of International Relations become
shorthand for entire approaches. In the process, qualifications and
subtleties are often lost. These iconic texts gain visibility through
citations and teaching, while younger scholars learn that professional
recognition may come from taking a distinctive sectarian position. New
``isms'' and specialized sections of professional associations can also
become vehicles for status and intellectual entrepreneurship. The result
is a centrifugal process in which extremism generates more extremism,
even if the original authors were more nuanced than later portrayals
suggest.

\subsection{The third pathology: confusing traditions with
theories}

The third pathology is the tendency to mistake a research tradition for
a complete theory. A tradition usually supplies broad assumptions, but
those assumptions are rarely sufficient to generate unique, deductively
valid predictions. Lake uses neorealism as an example. Waltz claimed
that anarchy and the desire for survival imply balancing behavior, yet
the same assumptions can also be compatible with cooperation, collective
security, or bandwagoning unless additional assumptions are introduced.
Similar indeterminacy appears in debates over relative versus absolute
gains in neorealism and neoliberalism. This is why several distinct
theories can coexist inside a single tradition---for example, offensive
realism, defensive realism, and neoclassical realism. They share some
core premises but differ in auxiliary assumptions. Because traditions
are often underspecified, they cannot be meaningfully ``tested'' as
wholes. Yet scholars frequently stage empirical contests between
traditions as if each generated clear predictions, producing debates in
which each side can interpret evidence as confirmation.

\subsection{The fourth pathology: selective empirical
focus}

The fourth pathology is empirical selectivity. Researchers tend to study
the topics, periods, and cases in which their preferred tradition is
most likely to perform well. Realists often focus on great-power
security, liberals on economic policy, institutionalists on
institutions, constructivists on norm change, and English School
scholars on socialization and international society. This pattern need
not result from deliberate bias; it can emerge naturally because
scholars test theories where they initially expect them to work, while
journals and publishers also favor positive findings. Nevertheless, the
cumulative consequence is self-confirmation. Each tradition studies its
strongest domain, finds supportive evidence, and then treats that
evidence as proof of general explanatory power. The research tradition
begins to function like theology: evidence is selected to affirm a prior
commitment rather than used to subject theory to demanding attempts at
falsification.

\subsection{The fifth pathology: the search for intellectual
hegemony}

The fifth pathology is the aspiration of each research tradition to
become the dominant scientific paradigm. Scholars often identify
themselves as realists, constructivists, or institutionalists in ways
that imply their preferred approach is broadly superior to alternatives.
Professional incentives reinforce this competition. If a tradition
becomes dominant, its leading scholars gain status, influence, and
control over disciplinary agendas. Because every group fears losing
ground to its rivals, competition for hegemony resembles an arms race:
even scholars who would prefer cooperation feel pressure to defend and
expand their approach. The collective result is worse for everyone.
Rather than acknowledging the limited scope of existing theories,
research communities expend enormous energy debating assumptions such as
whether actors are rational, whether structure or agency matters more,
or whether power, wealth, norms, or identity should be treated as
fundamental. Without comparable theories applied to the same phenomena,
these ``first-principle'' debates cannot be decisively resolved.

\subsection{Lake's Alternative: Problems, Partial Theories, and a Common
Lexicon}

Lake's preferred alternative is to reorganize scholarship around
substantive problems rather than around intellectual sects. Courses,
conference sections, and research agendas could focus on issues such as
war, climate change, development, inequality, genocide, or political
violence. The guiding question would then be ``What do we know about
this problem?'' rather than ``What are the core assumptions of this
school?'' Such an organization would encourage scholars to compare
explanations on the basis of their contribution to understanding a
phenomenon instead of their loyalty to a tradition.

A central part of this alternative is what Lake calls an ``embrace of
partiality.'' Theories should state their scope and boundary conditions
explicitly and identify what they do not explain. A theory is not
weakened merely because it applies only under certain conditions. On the
contrary, clarity about limits improves scientific evaluation because
scholars can distinguish genuine failures from cases outside a theory's
intended domain. Lake proposes that authors should routinely include a
brief statement of boundary conditions and explanatory limitations. This
would also make it easier to combine knowledge from different theories.
His metaphor is the familiar image of several people touching different
parts of an elephant: each theory may capture one part accurately, while
no single theory describes the entire animal.

Lake explicitly rejects atheoretical research. Theories remain necessary
for explaining patterns and identifying causes. His argument is instead
that scholars should use theories flexibly. A researcher might employ a
rational unitary-state model to study war, a theory of individual norms
to investigate global environmental change, and a sectoral
material-interest model to explain trade policy. Using different
theories for different problems should not be treated as intellectual
inconsistency if each theory is appropriate to the question at hand.

\subsection{Interests, interactions, and institutions as a ``Rosetta
stone''}

Analytic eclecticism creates a possible communication problem: if
scholars abandon fixed traditions, the field could become a ``tower of
Babel'' in which theories are too different to compare. Lake therefore
proposes a common lexicon built around three broad concepts: interests,
interactions, and institutions. Any complete theory of politics must
specify who the relevant actors are, what they want, how they interact,
and the institutional environment in which interaction occurs. Actors
may be individuals, groups, organizations, or states. Interests may
involve power, wealth, identity, status, security, or normative goals.
Interactions may involve zero-sum bargaining, cooperation,
socialization, or the reproduction of inequality. Institutions may range
from weak or epiphenomenal arrangements to dense structures that
constitute actors and enable communication itself.

This lexicon is deliberately neutral among traditions. A rationalist
theory may assume groups pursue material interests, while a
constructivist theory may explain interests as socially constituted. A
realist analysis may treat states as the principal actors, while another
approach disaggregates states into domestic groups. Yet both can still
be translated into the same basic questions about actors, interests,
interactions, and institutions. Components can also be modular:
something treated as given in one theory can be explained in another.
Regime type, for example, may be used as an independent variable in a
theory of war and then become the dependent variable in a separate
theory of democratization. The goal is not theoretical unification, but
clearer comparison, better specification, and cumulative knowledge.

\subsection{Epistemology: The Enduring
Divide}

Lake then turns from theory to epistemology. He identifies a deep divide
between nomological and narrative forms of explanation. This cleavage
has shaped major debates in international relations for decades and, in
his view, is harder to bridge than differences among theoretical
traditions. However, he does not believe that one epistemology will or
should defeat the other. Both have long histories, both respond to
different intellectual intuitions, and both can generate valuable
knowledge. The discipline should therefore accept epistemological
diversity rather than politicize it.

\subsection{The nomological approach}

The nomological approach seeks generalizable explanations based on
covering laws or theoretically specified relationships. It typically
uses the hypothetical-deductive method: researchers derive falsifiable
propositions of the form ``if X, then Y'' and test them using designs
that approximate controlled experiments. True experiments offer the
strongest internal validity, followed by quasi-experiments and then
correlational studies such as regression analysis. Causal claims can
emerge either from a strong design that isolates the effect of X on Y or
from a combination of theory and robust empirical association. Lake
notes that much international-relations research relies on the weaker
correlational form. Science in this tradition is iterative: theory
generates hypotheses, evidence tests them, and theories are revised,
extended, or occasionally rejected.

For Lake, generalizability in nomological research depends primarily on
theory rather than on the particular sample used in a test. If a theory
is explicitly intended to apply to all humans, evidence from a narrower
sample can still be relevant provided the test has strong internal
validity. Conversely, findings from a particular population cannot
validate a theory whose assumptions do not extend to that population.
External validity therefore rests on clearly specified scope and
boundary conditions.

\subsection{The narrative approach}

The narrative protocol organizes explanation differently. It begins with
a descriptive order that links events through time and examines how each
development shapes what follows. It then constructs a configurative
order by adjusting an interpretive framework to the historical facts
through a process of abduction. The objective is to produce a convincing
and plausible account of why events unfolded as they did. Narrative
explanations are often richer and more holistic than nomological models
because they emphasize combinations of causes and historically specific
conjunctures. Their weakness, in Lake's view, is that the criteria for
judging causal adequacy are more subjective. Narratives are evaluated in
terms of verisimilitude and interpretive credibility rather than
experimental internal validity.

Lake stresses that epistemology should not be confused with methodology.
A nomological study can use historical case studies and process tracing,
while a narrative analysis can employ large numbers of observations or
statistical techniques. ``Mixed methods'' can improve research by
testing theories in multiple ways, but mixing methods does not by itself
eliminate the distinction between forms of explanation. Method concerns
the techniques used to gather and analyze evidence; epistemology
concerns what counts as a persuasive explanation.

\subsection{Epistemological Sectarianism and Lake's
Recommendations}

The same pathologies that affect theory also affect epistemology.
Scholars reify the nomological and narrative approaches, reward extreme
positions, apply them selectively to problems where they are likely to
succeed, and then claim universal superiority. These preferences have
become politicized within the profession. Epistemological gatekeepers
can influence hiring, resource allocation, publication, and access to
professional networks, encouraging scholars to cluster with others who
share the same standards of explanation. Lake sees no objective basis
for declaring one epistemology universally superior. Scholars may simply
find one form of explanation more intuitively satisfying than another,
much as some people find beauty in mathematics and others in poetry.

Because neither epistemological tradition is likely to eliminate the
other, Lake proposes three principles for coexistence. First, scholars
should recognize the legitimacy of alternative forms of explanation and
respect colleagues who use them. Treating one's own epistemological
intuition as inherently superior is a form of intellectual hubris.
Second, each approach should be judged on its own terms. Narrative work
should be assessed for completeness, plausibility, competing
interpretations, and possible distortion of historical evidence.
Nomological work should be evaluated for logical deduction, valid
operationalization, and sound research design. Third, scholars should
accept that different approaches are likely to work better for different
questions.

Lake illustrates this final point with substantive examples. Narrative
approaches may be especially effective for slow, historically unique
changes in global norms, such as transformations in attitudes toward
slavery or human rights. Such processes unfold over long periods and are
difficult to analyze while holding the social environment constant. By
contrast, nomological approaches may be especially powerful when
interactions are frequent and institutions are stable enough to generate
repeated observations. Research on war and civil war, or on economic
policymaking within established democratic institutions, can often
benefit from this kind of analysis. The crucial question is therefore
not ``Which epistemology is best?'' but ``Which epistemology provides
the greatest insight for this particular problem?''

\subsection{Conclusion}

Lake concludes by returning to the social purpose of scholarship.
Academic privilege is justified by the expectation that scholars will
expand understanding. International studies fails in this responsibility
when it devotes too much energy to theological wars among theoretical
and epistemological camps. Progress comes from asking important
questions about phenomena that remain poorly understood, choosing
theories and methods appropriate to those questions, and evaluating
evidence honestly. At the present stage of the discipline, no
theoretical or epistemological approach merits hegemony. Diversity is
necessary, but diversity should not take the form of isolated
communities that speak mainly to themselves.

The essay's title is therefore deliberately provocative. ``Isms'' are
``evil'' not because realism, liberalism, constructivism, feminism, or
other traditions are useless, but because scholars too often transform
useful intellectual tools into identities and sects. Research traditions
become harmful when they are treated as complete theories, insulated by
selective evidence, and mobilized in struggles for disciplinary
dominance. Lake's alternative is a more modest and problem-centered
social science: partial theories with explicit limits, analytic
eclecticism across traditions, a shared vocabulary of interests,
interactions, and institutions, and mutual respect between nomological
and narrative forms of explanation. The ultimate measure of success
should be whether scholarship improves our understanding of world
politics and contributes, where possible, to social progress.
